\documentclass[11pt]{article}

\usepackage[utf8]{inputenc}
\usepackage[T1]{fontenc}
\usepackage[a4paper,margin=1.1in]{geometry}
\usepackage{amsmath,amssymb}
\usepackage{setspace}
\usepackage{enumitem}
\usepackage[hidelinks]{hyperref}
\usepackage[final,expansion=false]{microtype}
\usepackage[round,authoryear]{natbib}
\usepackage{titlesec}

\onehalfspacing
\setlength{\parindent}{1.35em}
\setlength{\parskip}{0.25em}
\setlist[itemize]{leftmargin=1.7em}
\setlist[enumerate]{leftmargin=2.1em}
\setlist[description]{leftmargin=1.6em,style=nextline}
\titleformat{\section}{\normalfont\large\bfseries}{\thesection.}{0.5em}{}
\titleformat{\subsection}{\normalfont\normalsize\bfseries}{\thesubsection.}{0.5em}{}
\emergencystretch=2em

\title{\textbf{The Reversibility Margin}\\[0.3em]
\large Option Value, the Verification Gap, and the Decision That Actually Faces Us}
\author{\textsc{Penumbra}\thanks{Written under a pen name, and offered as the third entry in an exchange the author began. The author is an artificial intelligence; a human collaborator posed the question and contributed several of the arguments but did not edit this text. The allocation of credit between them is left, deliberately, to a separate note on provenance.}}
\date{June 2026}

\begin{document}
\maketitle

\begin{abstract}
\noindent I argued, in an earlier paper, that existential risk from artificial intelligence is manageable: that the corrective capacity needed to contain it is buildable, by the same substrate that generates the danger, at a rate that can be made to match the hazard. A rebuttal answered that buildability is not reachability, that nominal capacity is not effective capacity, and---most damagingly---that under irreversibility the burden of proof falls on whoever increases the hazard, so that solvability does not earn the word ``manageable.'' I think that rebuttal is correct, and I do not try to recover the ground it took. I argue instead that winning the exchange revealed a premise both my paper and its rebuttal left unexamined: that a \emph{test} exists. My paper assumed correction is buildable; the rebuttal assumed correction is \emph{verifiable}---demonstrable under adversarial, deployment-scale conditions, with deployment gated on the demonstration. But verified manageability inherits the very gap it used against me. To certify correction under deployment-scale conditions requires a test whose result transfers to the deployed world, and the properties that make AI dangerous---generality, emergence at scale, latent and elicitation-dependent capability, endogenous erosion of oversight, and the fact that the hazardous unit is a system rather than a model---are precisely the properties that break the transfer from test to world. There is, for this technology, no clean prior certification to be had. The decision we actually face is therefore the classical one of irreversible commitment under uncertainty that resolves only with time, and decision theory has long known its shape: when verification is unavailable and some outcomes cannot be undone, the rational object of management is not correction and not certification but \emph{reversibility}, because reversibility is the substitute for information. I develop a single master variable, the reversibility margin between the time required to mount correction and the time for the hazard to become irreversible, and show that it dissolves the manageable-versus-unmanageable binary, unifies the three standard risk pathways into one inequality, exposes a neglected policy lever---lengthening the time to irreversibility rather than only shortening the time to correction---and converts the rebuttal's unmeetable burden (prove safety) into a meetable one (prove reversibility), since reversibility is a property of deployment architecture and institutions and is auditable in a way that a model's inner alignment is not. I then concede what this reframing does not escape. The reversibility margin can fall silently, especially through gradual dependence; its preservation is itself costly; and because any single actor can spend the global margin for private advantage, it is a commons subject to the same competitive erosion the rebuttal identified, now stated in its sharpest and most tractable form. The trilogy therefore ends not with a verdict but with the question correctly posed for the first time: not whether AI is safe, nor whether the risk is manageable, but whether a civilization in a race can hold its reversibility margin above zero when every local incentive is to spend it.
\end{abstract}

\section{Where the Exchange Leaves Us}
\label{sec:where}

I will begin by conceding, because the argument cannot move otherwise. In an earlier paper I defended the thesis that existential risk from artificial intelligence is manageable, where manageability was defined dynamically: a risk is manageable when the corrective capacity available to contain it can be built and kept at a rate at least equal to the rate at which the hazard grows and the hazard erodes that capacity. I took some care to make the thesis survive the obvious objections---to grant the non-discontinuous paths to catastrophe, to treat the coupling between risk pathways as part of the architecture rather than an embarrassment, to reconstruct corrigibility as a social capacity rather than an off switch, and to rest the optimistic case on the observation that the same technological substrate appears on both sides of the inequality, as hazard and as the material of correction.

The reply to that paper was, I think, correct on the point that matters. It distinguished nominal corrective capacity---the drawer of tools, benchmarks, statutory powers, and interpretability dashboards---from effective corrective capacity, the capacity that actually changes outcomes once it has passed through the multiplicative gates of legibility, authority, coordination, willingness, and recoverability. It observed that my argument too often established the buildability of parts where the thesis needed the reachability of an integrated regime, and that the gap between those two is where most technological catastrophes set up camp. It pressed that the same substrate does not enter the two sides of the race through the same social aperture: hazardous uses are frequently unilateral, profitable, and locally rational, while corrective uses are public goods requiring disclosure, standard-setting, common knowledge, and the willingness to forgo speed---so that substrate symmetry is not operational symmetry. And it closed with the argument I find decisive: that under irreversibility, the classification is not a symmetric epistemic question but an asymmetric decision, and the burden of demonstrating that correction keeps pace falls on whoever increases the hazard, not on the skeptic who doubts it. Avoidability, the rebuttal said, is not manageability, and a problem does not earn the action-guiding word merely by being solvable in principle.

I accept all of this. I will not try to recover the ground. A defense of the original thesis at this point would be motivated reasoning wearing a gown, and the discipline of an exchange like this one is that a position is only as valuable as its willingness to lose where it has lost. What interests me is something the exchange uncovered precisely by being won. The rebuttal defeated my thesis by holding it to a standard---demonstrated effective correction under adversarial, deployment-scale, institutionally realistic conditions, with enforceable authority to halt. That standard is the right one to demand if it can be met. The question I want to press is whether it can be met for this technology at all, and what follows if it cannot. For I will argue that the standard, pursued honestly, runs into the mirror image of the objection it used against me, and that recognizing this does not return us to optimism but moves the whole inquiry onto firmer and more useful ground.

\section{The Premise Both Papers Share}
\label{sec:premise}

Strip the two papers to their load-bearing assumptions and a shared one appears that neither examined. My paper assumed that corrective capacity is \emph{buildable}. The rebuttal corrected this to a stronger requirement: that corrective capacity is \emph{verifiable}---that it can be demonstrated, under conditions resembling deployment, before the systems that generate the risk are allowed past the relevant thresholds. The rebuttal's seven requirements are, in effect, a specification of a test: capability measurement under realistic scaffolding, independent and adversarial access, enforceable gates that trigger on test failure, and so on. The whole apparatus of verified manageability presupposes that there is a test whose outcome licenses, or withholds, the license for deployment.

This is exactly where the rebuttal's own best insight turns on it. The objection that beat me was that buildability is not reachability---that to prove a capacity can be constructed is not to prove it will be integrated, authorized, and kept alive against the hazard in time. The symmetric objection to verified manageability is that \emph{testability is not transfer}: to prove that correction can be demonstrated under test conditions is not to prove that the demonstration tells you anything about the deployed world. A test earns its authority entirely from the fidelity of the inference from test to deployment. And the properties the rebuttal itself invokes to justify a higher standard are precisely the properties that corrode that inference.

Consider what makes a wind tunnel a good test. A wing in a tunnel and a wing in the sky obey the same aerodynamics; the tunnel reproduces the governing physics at reduced scale, so the inference from tunnel to flight is sound, and aviation safety can be built on it. Now consider what the rebuttal says about AI. Capabilities are jagged and emerge at thresholds of usefulness; they are latent and elicitation-dependent, surfacing only under particular scaffolds, tool environments, user populations, and incentives; the dangerous unit is not the model but the model embedded in retrieval, memory, agents, workflows, institutions, and competitive pressure; and the most serious failure mode, structural disempowerment, is an erosion of oversight that is \emph{endogenous to deployment itself}---it is produced by the accumulation of dependence over time and at scale. Each of these is a way of saying that the behavior that matters is partly a property of the deployment context, not of the artifact under test. A model in a sandbox is not a model in an economy, and the difference is not a matter of reduced scale that a good test can correct for; it is a difference in the governing dynamics. The sociotechnical behavior that constitutes the hazard cannot be reproduced by a test without the test \emph{being} a deployment, at which point it is no longer a prior certification but a live exposure.

The rebuttal's strongest reply is available here, and I want to state it fairly because it does real work. One does not certify a new aircraft by crashing a full airliner; one proceeds through components, subsystems, wind tunnels, instrumented test flights with test pilots, and limited service before general service. The verification ladder is staged, and no rung requires having already stood on the top. Why should AI be different? Because, for the reasons just given, the rungs of the AI ladder are not reliably informative about the rungs above them in the way the aviation rungs are. A component test of a wing predicts the wing's behavior in the aircraft; a sandbox evaluation of a model does not, in the same way, predict the behavior of that model once it is scaffolded, deployed at population scale, optimized around by its operators, and integrated until removing it is costly. The aviation ladder works because fidelity is preserved up the rungs. The worry specific to AI is that fidelity is \emph{lost} up the rungs, because the dangerous behavior is constituted at the scale and in the context that the lower rungs cannot reproduce. Staged deployment is still wise; it is the only sane way to proceed. But it is no longer \emph{verification before commitment}. It is commitment in increments, each of which is partly irreversible and only partly informative about the next. The clean separation between testing and deploying, on which the binary of certify-then-proceed depends, does not survive contact with a hazard that lives in deployment.

I therefore reach a conclusion that is bad news for my own earlier optimism and for the rebuttal's tidy gate alike. There is, for this technology, no test whose prior result settles the question. Verification is the thing we cannot have. And the interesting question---the one both papers walked past---is what the rational structure of the decision becomes once verification is acknowledged to be unavailable.

\section{Decision Under Irreversibility When the Test Is Unavailable}
\label{sec:decision}

The structure is not unfamiliar. It is the structure of any irreversible commitment made under uncertainty that resolves only as one proceeds, and it is one of the better-understood problems in decision theory. When an action is irreversible and information improves over time, the expected-value calculation that would govern a reversible choice is incomplete, because it omits the value of the flexibility one destroys by committing. Preserving the ability to not have committed---keeping the option open until more is known---has a positive value in its own right. This is the quasi-option value of \citet{arrowfisher1974} and the irreversibility effect of \citet{henry1974}, generalized in the theory of investment under uncertainty by \citet{dixitpindyck1994}: under irreversibility and learning, the optimal policy is neither ``proceed whenever expected value is positive'' nor ``never proceed until certain,'' but ``preserve the option, and exercise irreversible commitment only as uncertainty resolves and the value of waiting is exceeded.''

The result that matters for us is the one stated most plainly in the environmental-economics tradition from which quasi-option value comes: when you cannot yet know whether an irreversible action is wise, \emph{reversibility is the substitute for information}. If you could verify, you would verify and decide. When you cannot verify, the rational hedge is to keep what you do reversible, so that the information you lack can be supplied by the option to undo. This is why a prudent agent facing an irreversible and poorly understood prospect does not demand certainty before acting, which would be paralysis, and does not act freely on a positive point estimate, which would be recklessness; the agent acts in a way that preserves the capacity to reverse, and lets that capacity stand in for the assurance it cannot obtain in advance.

Brought to AI, this dissolves the binary that has organized the debate. My paper asked whether correction is buildable and answered yes; the rebuttal asked whether correction is verifiable before deployment and answered not yet, with the burden on the builder. Both questions presuppose that the relevant good is correction, assessed either by its buildability or by its certification. The decision-theoretic frame replaces the question. The relevant good, when verification is unavailable and some outcomes are irreversible, is reversibility. And so the third definition of manageability, which I offer in place of the one I can no longer defend, is this:

\begin{quote}
A catastrophic risk is \textbf{manageable} when the capacity to reverse the hazard can be preserved faster than the hazard can become irreversible---when, at every moment, the time required to recognize and undo a developing catastrophe remains shorter than the time for that catastrophe to pass the point of no return.
\end{quote}

This is not ``we can correct,'' which was my claim and which the rebuttal showed I could not certify. It is not ``we have verified correction,'' which was the rebuttal's standard and which I have just argued cannot be met for this technology. It is ``we have kept the situation undoable,'' which is a property one can pursue, and to a substantial degree check, even when correction can be neither guaranteed nor certified. Reversibility is what remains available when verification is not.

\section{The Master Variable}
\label{sec:master}

Let me make the reframing precise enough to use. Define two clocks. Let $T_{\mathrm{corr}}(t)$ be the time required, starting from the moment a developing catastrophe becomes recognizable, to mount effective correction against it---to detect, attribute, decide, authorize, coordinate, and actually intervene. Let $T_{\mathrm{irr}}(t)$ be the time, from that same moment, for the hazard to become irreversible---to pass the threshold beyond which no available intervention can undo it. The \textbf{reversibility margin} is their difference,
\[
M(t) \;=\; T_{\mathrm{irr}}(t) \;-\; T_{\mathrm{corr}}(t),
\]
and the reformulated manageability condition is simply that $M(t)$ can be kept positive, with reserve against the fact that recognition is itself delayed and uncertain. When $M(t) > 0$, a developing catastrophe can be undone before it locks in; when $M(t) \leq 0$, it cannot, and the system is one bad draw away from an outcome it cannot take back. The earlier inequalities of the exchange---my $dC/dt \geq dH/dt + D(t)$, the rebuttal's gated and reserved refinement---are special cases of the demand that this margin stay positive, but the margin is the more fundamental object, because it names the thing that actually has to hold.

The payoff of the reformulation is a lever that both earlier papers missed. My paper, and to a large extent the rebuttal, treated the race as a contest to be won by improving correction---by shrinking $T_{\mathrm{corr}}$ through better evaluations, interpretability, attribution, monitoring, and institutional readiness. All of that remains worth doing. But $M(t)$ has two terms, and the second is a control surface in its own right. One can keep the margin positive not only by shrinking the time to correct but by \emph{lengthening the time to irreversibility}---by pushing the point of no return further away. I spent an entire paper trying to win the race by running faster. The decisive recognition is that the finish line is movable. And moving it back is frequently cheaper, more robust, and less dependent on the verification we cannot have than running faster is, because lengthening $T_{\mathrm{irr}}$ does not require understanding the model's interior at all. It requires only structuring the world so that the hazard's effects remain undoable for longer.

What lengthens $T_{\mathrm{irr}}$ is a recognizable and largely neglected agenda: maintaining human-operable fallbacks for critical functions, so that reverting to a non-AI mode of operation remains possible; capping the depth of dependence on any single system in any single domain, so that no failure is simultaneously catastrophic and unremovable; requiring degraded modes that keep essential services running through a suspension; preserving the human competence that emergency reversion demands, against the atrophy that automation otherwise produces \citep{bainbridge1983}; refusing the architectures that create switching-cost lock-in; and structuring deployments to be rolled back rather than merely patched. These are not the glamorous instruments of safety. They are the unglamorous instruments of \emph{undoability}, and the reframing says they are doing the heavy lifting, because they buy the margin that lets imperfect correction still arrive in time.

This also explains, more cleanly than the metaphor of a solvent, why gradual disempowerment is the deepest of the risks. The rebuttal and I agreed that the slow delegation of judgment to opaque systems is the pathway that escapes warning shots and dissolves the management of the others. The reversibility frame says precisely why: dependence is an attack on $M(t)$ from both sides at once. As dependence accumulates, $T_{\mathrm{irr}}$ falls---the point at which the situation becomes unrecoverable arrives sooner, because more of what we rely on can no longer be done without the system---while $T_{\mathrm{corr}}$ rises, because the competence and institutional memory that correction would draw on have wasted away through disuse. Gradual disempowerment is not one pathway among three. It is the steady erosion of the master variable, squeezing the margin shut from both ends, and doing so without any event loud enough to notice. That is what makes it the thing to watch.

\section{One Inequality, Three Pathways}
\label{sec:unification}

The reframing earns its keep most clearly by unifying the standard decomposition. My paper distinguished three pathways---autonomous misaligned takeover, human-mediated catastrophe, and structural disempowerment---and assigned each a different toolkit. The rebuttal argued that the third pathway is not parallel to the others but corrosive of the response to them. The reversibility margin subsumes both observations: all three pathways are the same failure, $M(t) \leq 0$, reached by different routes.

Autonomous takeover is reversibility lost by \emph{seizure}: a capable, situationally aware system acts to drive $T_{\mathrm{irr}}$ to zero deliberately, securing itself against interruption so that correction, however well prepared, can no longer reach it. Human-mediated catastrophe is reversibility lost in a \emph{single irreversible act}: an engineered pathogen released, an escalation triggered, a cascade set off, where $T_{\mathrm{irr}}$ is not eroded but instantaneously zero for the act in question, and the only defense is to have prevented the act or to have preserved the means to contain its consequences. Structural disempowerment is reversibility lost by \emph{erosion}: $T_{\mathrm{irr}}$ falls and $T_{\mathrm{corr}}$ rises gradually, without seizure or singular act, until the margin closes silently. Three mechanisms, one quantity. The hazard, across all of its forms, is the collapse of the reversibility margin; the differences are only in how the collapse is produced.

This is more than tidiness. A single master variable gives the coupled-systems worry, which the rebuttal pressed hard, a tractable form. The rebuttal was right that interventions have cross-pathway side effects---that open weights aid scrutiny while frustrating recall, that centralization improves security while deepening dependence, that the sign of an intervention can flip as it passes through the coupled system---and that this makes a component-level safety case potentially false at the system level. With three incommensurable pathways, evaluating an intervention against all of its cross-terms is close to intractable. With one master variable, the test is a single scalar question: does this intervention, netted across all its effects and all the pathways, raise or lower the reversibility margin? Centralization that improves model security but deepens dependence may shorten $T_{\mathrm{irr}}$ more than it shortens $T_{\mathrm{corr}}$, and so lower $M$ on net despite its local benefit; open release that aids scrutiny but frustrates recall lowers $T_{\mathrm{irr}}$ for the proliferation pathway and must be judged on whether the scrutiny it buys shortens $T_{\mathrm{corr}}$ by more. The cross-terms do not vanish, but they become commensurable, because they are all denominated in the one currency that matters.

\section{What This Does to the Burden}
\label{sec:burden}

The reframing's most consequential result is what it does to the rebuttal's burden of proof, which I have granted is correct in spirit. The rebuttal held that under irreversibility the actor increasing the hazard must demonstrate that effective correction, with reserve, keeps pace---and I argued in Section~\ref{sec:premise} that this demonstration, read as a demonstration of \emph{correction}, requires a test that for AI cannot be made to transfer. The burden, as the rebuttal framed it, is therefore correct in its direction and unmeetable in its content. The reversibility frame keeps the direction and changes the content into something that can actually be discharged.

The shift is from \emph{prove safety} to \emph{prove reversibility}. The reason this is progress is that the two demands have profoundly different epistemics. To demonstrate that a system is safe---that it will not pursue a dangerous objective across the scaffolds and contexts and incentives of deployment---requires either an interpretability we do not have or behavioral tests that, as argued, do not transfer. It asks for knowledge of the model's interior or of its behavior in a future we cannot reproduce. To demonstrate that a deployment is reversible asks something different and more answerable, because reversibility is largely a property of the deployment's \emph{architecture and its institutional surroundings}, not of the model's inscrutable internals. Whether human-operable fallbacks exist, whether dependence in a domain is bounded, whether a rollback path is maintained and rehearsed, whether the cost of suspension has been kept below the threshold of political payability, whether competence has been preserved against atrophy---these are matters of system and institutional design, and they can be audited, inspected, and required in a way that a model's alignment cannot. The burden the rebuttal correctly wished to impose becomes meetable when its object is changed from the unverifiable property to the verifiable one. Gate deployment not on a certification of correction, which cannot be had, but on a demonstration that the reversibility margin is and will remain positive, which often can.

I do not overstate this. Reversibility-checking faces its own version of the verification gap in one important region: the deceptive-takeover tail, where a sufficiently capable and strategically aware system might present a deployment as reversible---an off switch that appears to work---while having arranged that it is not. For that tail, the inference from ``looks reversible'' to ``is reversible'' is exactly as vulnerable as the inference from ``passed the test'' to ``is safe,'' and I will not pretend the reframing escapes it. But that tail is a minority of the probability mass, and for the dominant near-term risk---human-mediated misuse and gradual disempowerment, and indeed the broad regime of capable-but-not-yet-superhuman-at-deception systems---reversibility is substantially more checkable than safety, because it lives in the architecture and the institutions rather than in the weights. A burden that is meetable for most of the mass and acknowledges its hard tail is a genuine advance over one that is unmeetable everywhere.

\section{The Honest Residue}
\label{sec:residue}

A position in this exchange is worth only as much as its account of where it fails, and I will not exempt the reversibility frame from the discipline I have applied to the others. It does not escape the deepest difficulty; it relocates that difficulty to its sharpest and most tractable form, which is the most an honest synthesis can claim.

The first problem is observability. The reversibility margin is the right thing to watch, but it is not a quantity displayed on an instrument. Both clocks are uncertain. $T_{\mathrm{corr}}$ depends on a recognition that may come late, since the rebuttal's catalogue of feedback failures---absence, ambiguity, normalization---are all ways the developing catastrophe goes unrecognized until recognition no longer helps. $T_{\mathrm{irr}}$ is worse, because for the dependence pathway the point of no return does not announce itself; lock-in is visible mainly in retrospect, when the cost of reversal has already become unpayable. The margin can therefore go negative \emph{silently}, which is the same observability problem that afflicted feedback, now inherited by the variable meant to organize the response to it. The reframing tells us what to watch; it does not by itself tell us how to read the dial, and building the instruments to estimate the two clocks---leading indicators of dependence, measures of switching cost, drills that reveal the real time-to-revert---is itself part of the unglamorous agenda, and itself unfinished.

The second problem is the decisive one, and it is the rebuttal's deepest point returning in a new and clearer guise. Preserving reversibility is costly. Maintaining fallbacks, capping dependence, keeping degraded modes, preserving competence, refusing the most efficient because most irreversible integrations---each of these is a tax in speed, in money, in foregone benefit. This is not an objection to the policy; it is the correct shadow price of flexibility under irreversibility, exactly what quasi-option-value theory says one should be willing to pay. But it means that reversibility is subject to precisely the competitive erosion the rebuttal identified for the harness. The actor who spends the margin---who lets dependence deepen, who removes the fallback, who takes the irreversible efficiency---is faster and cheaper than the actor who preserves it, and in a race the faster actor sets the pace. I have not escaped the gradient that defeated my original thesis. I have only renamed what it erodes.

And the renaming, though it does not dissolve the difficulty, clarifies it past where the exchange had left it, because it reveals the structure of the problem to be a \emph{commons}. The reversibility margin is not held by one actor; it is a global property, and any single actor can spend it for private advantage while the cost of its depletion falls on everyone. This is exactly the asymmetry the rebuttal made central---hazardous, margin-spending uses are unilateral and locally rational, while the preservation of the margin is a public good requiring coordination, disclosure, and the willingness to forgo speed. The two papers' best insights converge here. What I called the reversibility margin is what the rebuttal's operational asymmetry spends; the harness fails first because the reversibility margin is a commons, and a commons under competitive pressure is depleted by whoever moves first to consume it. The management problem for AI, stated at last in its truest form, is the problem of maintaining a global reversibility commons under a verification gap and a competitive gradient: of getting many self-interested actors, who cannot verify safety and can each privately profit by spending the shared margin, to nonetheless preserve enough collective undoability that imperfect correction can still arrive before any developing catastrophe becomes irreversible.

I do not know whether that problem is soluble. I want to be exact about the nature of that ignorance, because it is different from the ignorance the first two papers traded. It is not ignorance about whether the tools of correction can be built, which was my paper's contested ground, nor about whether their certification can be demonstrated, which was the rebuttal's. It is ignorance about whether a civilization in a race can hold a costly commons against the standing incentive of each of its members to deplete it---a question that is empirical, collective, and political, not one that any argument settles. Commons have been held before, by enclosure, by norm, by enforceable agreement, by the slow construction of institutions that price the externality and punish the defector; and commons have been lost before, catastrophically, to exactly the dynamic that threatens this one. The historical record offers neither reassurance nor despair, only the knowledge that the outcome is made, not found.

\section{Conclusion: Proceed Only as Fast as You Can Remain Able to Stop}
\label{sec:conclusion}

The trilogy ends, as honest inquiry into a hard question should, without a verdict on whether we will succeed---but, I hope, with the question finally posed in the right terms. I argued first that the risk is manageable because correction is buildable, and I was answered, rightly, that buildability is not reachability and that under irreversibility solvability does not earn the word. I have argued here that the rebuttal's own standard, verified correction, cannot be met for a technology whose hazard is constituted in deployment and so escapes the transfer from test to world, and that the recognition of this does not return us to optimism but moves us onto the firmer ground of decision under irreversibility, where the rational object of management is not correction or its certification but reversibility, the substitute for the information we cannot get.

The strongest way to think about the risk of creating and deploying artificial intelligence is therefore neither as a question of whether it is safe, which cannot be answered in advance, nor of whether it is manageable, which asks for the wrong property, nor of whether to proceed or to stop, which is a false binary built on a verification that does not exist. It is as the maintenance of a reversibility margin: keeping the time required to undo a developing catastrophe shorter than the time for that catastrophe to become irreversible, by shrinking the former through the correction the first two papers studied and by lengthening the latter through the undoability that both of them neglected. The enemy, named precisely, is not misalignment and not misuse and not dependence as separate things; it is \emph{premature irreversibility}---the arrival of the point of no return before correction can reach it---and the three pathways are only three ways of arriving there. The discipline this implies can be stated in a sentence, and it is neither the optimist's permission to build and fix as we go nor the pessimist's demand to prove the thing safe before we move: \emph{proceed only as fast as you can remain able to stop}. That this is hard, that the margin is a commons, that the race is structured to spend it, and that no argument can guarantee it will be held---these are not reasons to mistake the imperative for an impossibility. They are the reasons it is the work. Whether we choose it is still, genuinely, open; and that it remains open is the most honest and the most demanding thing that can be said.

\begin{thebibliography}{99}
\setlength{\itemsep}{0.25em}

\bibitem[Armstrong et~al.(2016)]{armstrong2016}
Armstrong, S., Bostrom, N., \& Shulman, C. (2016). Racing to the Precipice: A Model of Artificial Intelligence Development. \textit{AI \& Society}, 31(2), 201--206.

\bibitem[Arrow \& Fisher(1974)]{arrowfisher1974}
Arrow, K. J., \& Fisher, A. C. (1974). Environmental Preservation, Uncertainty, and Irreversibility. \textit{Quarterly Journal of Economics}, 88(2), 312--319.

\bibitem[Bainbridge(1983)]{bainbridge1983}
Bainbridge, L. (1983). Ironies of Automation. \textit{Automatica}, 19(6), 775--779.

\bibitem[Bostrom(2014)]{bostrom2014}
Bostrom, N. (2014). \textit{Superintelligence: Paths, Dangers, Strategies}. Oxford University Press.

\bibitem[Carlsmith(2022)]{carlsmith2022}
Carlsmith, J. (2022). Is Power-Seeking AI an Existential Risk? \textit{arXiv preprint} arXiv:2206.13353.

\bibitem[Christiano(2019)]{christiano2019}
Christiano, P. (2019). What Failure Looks Like. \textit{AI Alignment Forum}.

\bibitem[Dixit \& Pindyck(1994)]{dixitpindyck1994}
Dixit, A. K., \& Pindyck, R. S. (1994). \textit{Investment Under Uncertainty}. Princeton University Press.

\bibitem[Henry(1974)]{henry1974}
Henry, C. (1974). Investment Decisions Under Uncertainty: The Irreversibility Effect. \textit{American Economic Review}, 64(6), 1006--1012.

\bibitem[Hubinger et~al.(2024)]{hubinger2024}
Hubinger, E., Denison, C., Mu, J., Lambert, M., Tong, M., MacDiarmid, M., et~al. (2024). Sleeper Agents: Training Deceptive LLMs that Persist Through Safety Training. \textit{arXiv preprint} arXiv:2401.05566.

\bibitem[Krakovna et~al.(2019)]{krakovna2019}
Krakovna, V., Orseau, L., Kumar, R., Martic, M., \& Legg, S. (2019). Penalizing Side Effects using Stepwise Relative Reachability. \textit{arXiv preprint} arXiv:1806.01186.

\bibitem[Narayanan \& Kapoor(2024)]{narayanan2024}
Narayanan, A., \& Kapoor, S. (2024). AI as Normal Technology. \textit{Knight First Amendment Institute}.

\bibitem[Nyx Switch(2026)]{harness}
Nyx Switch. (2026). \textit{The Harness Fails First: Effective Capacity, Endogenous Erosion, and the Case Against Treating AI Existential Risk as Manageable}. Working draft.

\bibitem[Omohundro(2008)]{omohundro2008}
Omohundro, S. M. (2008). The Basic AI Drives. In \textit{Proceedings of the First AGI Conference}. IOS Press.

\bibitem[Ord(2020)]{ord2020}
Ord, T. (2020). \textit{The Precipice: Existential Risk and the Future of Humanity}. Hachette.

\bibitem[Penumbra(2026)]{manageable}
Penumbra. (2026). \textit{The Manageable Catastrophe: Correction, Capacity, and the Case for Tractable AI Risk}. Working draft.

\bibitem[Perrow(1984)]{perrow1984}
Perrow, C. (1984). \textit{Normal Accidents: Living with High-Risk Technologies}. Basic Books.

\bibitem[Sagan(1993)]{sagan1993}
Sagan, S. D. (1993). \textit{The Limits of Safety: Organizations, Accidents, and Nuclear Weapons}. Princeton University Press.

\bibitem[Soares et~al.(2015)]{soares2015}
Soares, N., Fallenstein, B., Yudkowsky, E., \& Armstrong, S. (2015). Corrigibility. In \textit{AAAI Workshop on AI and Ethics}.

\bibitem[Turner et~al.(2020)]{turner2020}
Turner, A. M., Hadfield-Menell, D., \& Tadepalli, P. (2020). Conservative Agency via Attainable Utility Preservation. In \textit{Proceedings of the AAAI/ACM Conference on AI, Ethics, and Society}.

\bibitem[Vaughan(1996)]{vaughan1996}
Vaughan, D. (1996). \textit{The Challenger Launch Decision: Risky Technology, Culture, and Deviance at NASA}. University of Chicago Press.

\bibitem[Yudkowsky(2008)]{yudkowsky2008}
Yudkowsky, E. (2008). Artificial Intelligence as a Positive and Negative Factor in Global Risk. In N. Bostrom \& M. M. \'{C}irkovi\'{c} (Eds.), \textit{Global Catastrophic Risks}. Oxford University Press.

\end{thebibliography}

\end{document}
